\chapter{Problem Statement}
\label{chap:problem-statement}

\section{Federated Learning Setting}
\label{sec:fl-setting}

FL trains a shared model across $K$ clients, where client $k$ keeps a
private dataset $\mathcal{D}_k$ and the server optimizes the global
objective
\begin{equation}
\min_{w} \; F(w) = \sum_{k=1}^{K} \frac{n_k}{n} \, F_k(w),
\qquad
F_k(w) = \frac{1}{n_k} \sum_{i \in \mathcal{D}_k} \ell(w; x_i, y_i),
\label{eq:fl-objective}
\end{equation}
where $w$ are the shared parameters and $n_k = |\mathcal{D}_k|$, $n = \sum_k n_k$.
In round $t$, the server sends $w_t$ to a subset of clients; each client
trains locally and returns an update $\Delta w_k^{(t)}$ (or, in the
single-step case this project focuses on, a gradient
$g_k^{(t)} = \nabla_w F_k(w_t)$ scaled by the learning rate). The server
aggregates with the FedAvg rule
\begin{equation}
w_{t+1} = w_t + \sum_{k \in S_t} \frac{n_k}{\sum_{j \in S_t} n_j} \, \Delta w_k^{(t)}.
\label{eq:fedavg}
\end{equation}
Only $\Delta w_k^{(t)}$ (equivalently $g_k^{(t)}$) ever leaves the client ---
never $x_i$ or $y_i$ directly.

\section{Threat Model}
\label{sec:threat-model}

We consider an \emph{honest-but-curious} attacker: it follows the FL
protocol exactly as specified, but additionally tries to extract private
data from whatever it legitimately observes. Two placements are relevant,
and both reduce to the same observation:

\begin{itemize}[leftmargin=1.6em]
  \item \textbf{A curious server}, which by design receives every client's
  update $\Delta w_k^{(t)}$ (or gradient $g_k^{(t)}$) each round.
  \item \textbf{An eavesdropper} on the client--server channel, which
  intercepts the same quantity in transit.
\end{itemize}

\textbf{What the attacker knows:} the model architecture and loss function
$\ell(\cdot\,;\cdot)$, the current global weights $w_t$, and the leaked
gradient (or update) itself. \textbf{What the attacker does not know:} the
private batch $\{(x_i, y_i)\}$ that produced that gradient. The attacker
makes no other assumption and does not deviate from the protocol --- it is
purely a passive observer of data every FL deployment already produces.

\section{The Gradient Inversion Problem}
\label{sec:gi-problem}

Formally, for a private batch of size $B$,
\begin{equation}
g = \nabla_w \, \ell\big(w_t;\, \{x_i, y_i\}_{i=1}^{B}\big)
\label{eq:leak}
\end{equation}
is the quantity the attacker observes (for $B=1$ this is simply
$\nabla_w \ell(w_t; x, y)$). Equation~\eqref{eq:leak} is a deterministic
function of the private batch and the known weights $w_t$; it is not
necessarily one-way. The \textbf{gradient inversion problem} is to invert
it:

\begin{quote}
\textbf{Input:} the model architecture and loss $\ell$, the weights $w_t$,
and the leaked gradient $g$. \\
\textbf{Output:} an estimate $\{(\hat{x}_i, \hat{y}_i)\}_{i=1}^{B}$ of the
private batch.
\end{quote}

The two attack families surveyed in Chapter~\ref{chap:related-work} answer
this input/output contract differently. Optimization-based attacks pose it
as a search problem,
\begin{equation}
\min_{\hat{x}, \hat{y}} \; D\big(g', g\big), \qquad g' = \nabla_w \, \ell(w_t; \hat{x}, \hat{y}),
\label{eq:opt-attack}
\end{equation}
for some distance or similarity measure $D(\cdot,\cdot)$ between gradient
vectors, solved by gradient descent on $(\hat{x}, \hat{y})$. Closed-form
attacks instead treat Equation~\eqref{eq:leak} as a system of per-layer
equations and solve directly for $\{\hat{x}_i\}$ by linear algebra, with no
search and no loss $D$ to minimize. Chapter~\ref{chap:methodology} details
both routes for each method implemented in this project. How the quality
of $(\hat{x}, \hat{y})$ is measured against the true $(x, y)$ --- and under
what gradient-perturbation defense budget --- is defined separately in
Chapter~\ref{chap:experiment}, since it is a property of the evaluation, not of the problem
statement itself.
